\scanpage{7}\mspage{\handwritten{6}}

morphismes de flêches.

Notons une symmétrie évidente dans la catégorie $(\mathrm{Bicat})$, interprétée en termes d\deleted{e}'ensemble bi-semisimpliciaux, obtenue en échangeant les \unsure{lignes} et les colonnes dqns \deleted{les} de tels ensembles bi... Un ensemble \unsure{bisemisimplicial} qui ... peut donc de deux façons différentes être considéré comme définissant une bicatégorie. De même le groupe symétrique $\mathrm{Sym}_n$ opère sur la catégorie \deleted{de} $(n\unsure{-|{}}\mathrm{upcat})$ des catégories $n$-uples ..

Une 2-catégorie peut être définie comme \unsure{constttuée|constituée} par un ensemble \deleted{discret} $C^{\deleted{o}}_{\handwritten{0}}$ (l'ensemble des objets ), qu'on interprètera comme une catégorie discrète, plus la donnée pour tout $x,y\handwritten{\in}C^{\deleted{o}}_{\handwritten{0}}$ d'une catégorie $\underline{\mathrm{Hom}}(x,y)$, plus \deleted{une} \inserted{des foncteurs} composition entre celles-ci ... Prenant la catégorie somme \deleted{de} $C^{\deleted{1}}_{\handwritten{1}}$ des $\underline{\mathrm{Hom}}(x,y)$, On peut exprimer la donnée de ces foncteurs de composition, et des sources et buts \inserted{$x,y$} dans les catégories $\underline{\mathrm{Hom}}(x,y)$, comme définis par des foncteurs ``source'' et ``but'' $C^{\deleted{1}}_{\handwritten{1}}\handwritten{\to}C^{\deleted{o}}_{\handwritten{0}}$, et un foncteur de composition $C^{\deleted{1}}_{\handwritten{1}}\mathbin{x}_{C^{\deleted{o}}_{\handwritten{0}}}C^{\deleted{1}}_{\handwritten{1}}\handwritten{\to}C^{\deleted{1}}_{\handwritten{1}}$, les axiomes des 2-catégories correspondant exactement aux axiomes des$(\mathrm{Cat})$-catégories. De cette façon, la catégorie des \unsure{2-catégories} est équivalente à la catégorie des $(\mathrm{Cat})$-catégories qui ont une ``\unsure{catégorie} d'objets'' qui est discrète, ou encore à la catégorie des Bicatégories qui, interprétées comme ensembles bi-semisimpliciaux, ont \deleted{une} comme première ligne un ensemble \unsure{bemisimplicial|semisimplicial} discret (i.e. un foncteur constant sur $\handwritten{\Delta}$). Il parait qu'on a une interprétation simple dans le même esprit de la notion de $n$-catégorie en termes d'ensembles \unsure{$n$-semi-simpliciaux}.

\margin{\handwritten{\(\Vert\)}}% double marginal stroke beside Demander
\deleted{au} Demander à Quillen relations entre la notion de $m$-Catégorie dans la catégorie $(n\text{-}\mathrm{Cat})$, et celle d\deleted{e}\handwritten{'} ensemble $(m+n)$-semisimplicial. \deleted{(en particulier entre cette dernière notion et celle de}

L'interprétation standard semisimpliciale permet d'interpréter

% archive sheet number: 6
